\documentclass{article}
\usepackage{amsmath}
\usepackage{amssymb}
\usepackage{amsthm}
\usepackage[T1]{fontenc}


\newtheorem*{conjecture}{Conjecture}
\newtheorem*{theorem}{Theorem}
\newtheorem{lemma}{Lemma}
\newtheorem*{lemma*}{Lemma}
\newtheorem*{proposition}{Proposition}


\makeatletter
\newcommand*{\rom}[1]{\expandafter\@slowromancap\romannumeral #1@}
\makeatother


\begin{document}

\section*{Problem 9*}

First, observe that \( 0 < \frac{1}{n} \) for every \(n \geq 1\) and \(0\) is in every \(A_n\). Then \( 0 \in A \).
\newline
It is known that \(A \subseteq [0, \infty)\). Indeed, if \(a \in A\) then:
\begin{equation*}
    (\forall n \geq 1) \ \ a \in A_{n} \implies a \in A_{1} = [0, 1) \implies a \in [0, 1] \implies a \in [0, \infty)
\end{equation*}

Assume now that \(a > 0\). We know that \(\lim_{n \to \infty} \frac{1}{n} = 0\).
Let \(\epsilon = a\).  From the definition of the limit there exists \(n_0 \in \mathbb{N}\) such that:

\begin{equation*}
    a = \epsilon > |\frac{1}{n} - 0| = |\frac{1}{n}| = \frac{1}{n} \ \ \mbox{dla} \ \ n \geq n_0
\end{equation*}

Then:

\begin{equation*}
    a > \frac{1}{n} > 0 \ \ \mbox{dla} \ \ n \geq n_0
\end{equation*}

We have \(a \notin A_{n_0}\). Using contraposition:

\begin{equation*}
    \Big( a \notin \{0\} \implies a \notin A \Big) \iff \Big( a \in A \implies a \in \{0\} \Big)
\end{equation*}

and finally we have:

\begin{equation*}
    A = \bigcap_{n=1}^{\infty} [0, \frac{1}{n}) = \{ 0 \}
\end{equation*}

\end{document}
